\chapter{Visualization and Rendering of Semantic Galaxies}
\label{ch:visualization}

The preceding chapters developed a full dynamical system for
semantic galaxy morphogenesis:
RSVP fields evolve on the manifold (\mathcal{M}),
Polyxan atoms follow an N-body flow,
and sheaf-derived cohomological forces enforce global consistency.
This chapter establishes the pipeline for \emph{visualizing}
these structures, both for scientific analysis and for interactive
exploration of semantic geometry.

Visualization here is not merely a rendering task.
It is the final computational layer through which semantic coherence,
field structurization, galaxy morphogenesis, and higher cohomological
patterns become directly inspectable by a human or automated observer.
The aim is to produce:
\begin{itemize}
\item real-time trajectories of Polyxan atoms,
\item field slices and isosurfaces for (\Phi,\mathbf{v},S),
\item curvature, density, and sheaf obstruction maps,
\item global and local visual invariants,
\item interactive 3D navigation and diagnostic tools.
\end{itemize}

We treat visualization as a structured, multi-resolution sheaf over
rendering layers, ensuring correctness and coherence.


% ============================================================
\section{Principles of Semantic Visualization}
% ============================================================

Semantic galaxies occupy a high-dimensional manifold (\mathcal{M}),
but visualization is typically performed in 2D or 3D.
Thus a projection functor is required:
\[
  \Pi : \mathcal{M} \longrightarrow \mathbb{R}^d, \qquad d \in \{2,3\}.
\]

\begin{definition}[Semantic Projection Functor]
A map $\Pi$ is a semantic projection if:
\begin{enumerate}
  \item It preserves adjacency: 
        $d_{\mathcal{M}}(x,y) \ll 1 \Rightarrow 
         \|\Pi(x)-\Pi(y)\| \ll 1$.
  \item It respects RSVP gradient directions: 
        $\Pi_*(\nabla\Phi) \approx \text{principal flow directions}$.
  \item It pushes forward sheaf coherence:
        local sections remain visually compatible after projection.
\end{enumerate}
\end{definition}

Typical choices include:
\begin{itemize}
\item PCA or spectral embeddings on the manifold discretization,
\item diffusion maps weighted by $S$,
\item t-SNE or UMAP for interactive exploration,
\item intrinsic mesh flattening for 2D field slices.
\end{itemize}

The projection $\Pi$ becomes part of the visualization sheaf:
\[
  \mathscr{V}(U) = \text{renderable slices of } \Pi(U).
\]


% ============================================================
\section{Rendering RSVP Fields}
% ============================================================

The three RSVP fields $(\Phi,\mathbf{v},S)$ can be visualized directly.

\subsection{Scalar Field Rendering}

The scalar fields $\Phi$ and $S$ may be shown as:
\begin{itemize}
\item color maps on surfaces or point clouds,
\item height maps in 2.5D,
\item contour plots or isosurfaces.
\end{itemize}

Define the color operator:
\[
  \mathcal{C}_\Phi(x)
    = \mathrm{colormap}_\Phi\big(\Phi(x)\big),
\qquad
  \mathcal{C}_S(x)
    = \mathrm{colormap}_S\big(S(x)\big),
\]
with colormaps chosen to emphasize curvature and gradient magnitudes.

\subsection{Vector Field Rendering}

For the vector field $\mathbf{v}$:
\begin{itemize}
\item arrows or glyphs,
\item line integral convolution (LIC),
\item streamline animations,
\item tangent-field ribbons.
\end{itemize}

We define the streamline evolution:
\[
  \gamma'(t) = \mathbf{v}(\gamma(t)),\quad
  \gamma(0)=x_0,
\]
which forms the basis of dynamic visualizations
such as semantic characteristic curves.


% ============================================================
\section{Rendering Polyxan Atoms and Hyperedges}
% ============================================================

Each Polyxan atom $v$ appears as a point or glyph at $\Pi(x_v)$.
Atoms may be colored by:
\begin{itemize}
\item semantic mass $m_v$,
\item local curvature $H(v)$,
\item density defect $|\rho_\iota(v)-\rho^*(v)|$,
\item sheaf obstruction magnitude.
\end{itemize}

Hyperedges appear as:
\begin{itemize}
\item line segments (binary relations),
\item polygonal surfaces (ternary relations),
\item translucent “semantic membranes’’ (high-arity relations).
\end{itemize}

For high-degree hyperedges,
it is often preferable to decorate atoms with local frames indicating
their incidence structure rather than rendering explicit surfaces.


% ============================================================
\section{Visualization of Sheaf Cohomology}
% ============================================================

Cohomological obstructions are global objects that must be rendered locally.

We define the \emph{obstruction field}:
\[
  \mathfrak{O}(x_v)
    := \sum_{U_i\ni v}
       \gamma_{U_i}\|c^1_{i,j}\|
       + \gamma'_{i,j,k}\|c^2_{i,j,k}\|.
\]

Visualization strategies include:
\begin{itemize}
\item heatmaps over particle positions,
\item isosurfaces showing boundaries of semantic inconsistency,
\item blinking or pulsing glyphs at high-obstruction nodes,
\item topological summaries via persistence diagrams.
\end{itemize}

For global cohomology,
we display:
\[
  \dim H^0(\mathcal{G},\mathcal{F}),\quad
  \dim H^1(\mathcal{G},\mathcal{F}),\quad
  \dim H^2(\mathcal{G},\mathcal{F}),
\]
as a time-series
tracking semantic coherence over simulation steps.


% ============================================================
\section{Dynamic Visualization of Embedding Flows}
% ============================================================

To illustrate embedding motion,
we display time-dependent trajectories
\[
  t \mapsto \Pi(x_v(t)).
\]

Use:
\begin{itemize}
\item trail paths,
\item velocity arrows,
\item acceleration glyphs,
\item blending to show recent movement.
\end{itemize}

Particles may be grouped into semantic clusters detected via:
\[
  C_k = \{v : x_v \in \text{basin of attraction of local }\Phi\text{ minimum}\}.
\]

Clusters are rendered with shared color hues or bounding surfaces.


% ============================================================
\section{Isosurfaces and Field Slicing}
% ============================================================

For RSVP fields on a 3D mesh:

\begin{itemize}
\item choose a slicing plane $P$,
\item define $\Phi_P = \Phi|_P$, $S_P = S|_P$,
\item render level sets,
\item overlay Polyxan atoms intersecting $P$.
\end{itemize}

Alternatively, extract isosurfaces via marching cubes:
\[
  \Phi(x)=\Phi_0, \qquad S(x)=S_0.
\]

These show semantic valleys, entropy ridges, and flow bottlenecks.


% ============================================================
\section{Multi-Resolution Visualization Sheaf}
% ============================================================

Visualization is organized as a sheaf
\[
  \mathscr{V} : \mathrm{Open}(\mathcal{M}) \to \mathrm{Scenes}
\]
where $\mathscr{V}(U)$ is the scene graph for region $U$.

\subsection{Local Layers}

Each region has:
\begin{itemize}
\item field slices,
\item local atoms,
\item neighborhood hyperedges,
\item cohomology diagnostics.
\end{itemize}

\subsection{Global Layers}

The global scene is the colimit:
\[
  \mathrm{Scene}_{\mathrm{global}}
    = \mathrm{colim}_{U_i} \mathscr{V}(U_i).
\]

This ensures visual consistency:
objects agree on overlaps whenever sheaf conditions hold.

Derived functors produce higher-level visuals:
\begin{itemize}
\item $\mathbf{R}^1\Gamma$ visualized as loops of inconsistency,
\item $\mathbf{R}^2\Gamma$ visualized as voids or semantic cavities.
\end{itemize}


% ============================================================
\section{Interactive Navigation and Tooling}
% ============================================================

Real-time exploration is supported by:

\begin{itemize}
\item camera orbit around dense clusters,
\item semantic zooming via reparameterization of $\Pi$,
\item focus+context rendering,
\item selector tools for atoms or hyperedges,
\item temporal scrubbing across simulation history,
\item toggling sheaf layers and field layers.
\end{itemize}

Interaction must preserve projection coherence:
user actions should not violate adjacency or flow invariants.


% ============================================================
\section{Summary}
% ============================================================

This chapter established the visualization stack for semantic galaxies:

\begin{itemize}
  \item projection functors for mapping (\mathcal{M}) to view space,
  \item rendering of RSVP fields and Polyxan atoms,
  \item visualization of curvature, density, and cohomological invariants,
  \item dynamic embedding-flow trajectories,
  \item multi-resolution, sheaf-consistent scene construction,
  \item interactive navigation tools.
\end{itemize}

Visualization is thus itself a sheaf-theoretic layer,
ensuring semantic and geometric coherence even in reduced dimensions.

The next chapter will integrate control loops and user-defined forces,
allowing interactive manipulation and steering of semantic galaxies.

